\section{Cyber Resilience Act Quick-Scan}
\label{sec:cra}

% TODO: Cyber Resilience Act (CRA - Regulation (EU) 2024/2847) Quick-Scan for SecureBank PT
% 
% 1. Product Scope and Classification:
%    - Define product boundary: Mobile banking application (Android/iOS) distributed to consumers
%      qualifies as a "Product with Digital Elements" (PDE) under CRA Article 2.
%    - Delineate boundary with backend/cloud banking services covered under DORA (Regulation (EU) 2022/2554)
%      and NIS2 (Directive (EU) 2022/2555).
%    - CRA category assessment: Standard product with digital elements vs. Important/Critical product
%      (Annex III/IV) based on authentication and financial access role.
%
% 2. Applicable Essential Requirements (Annex I):
%    - Part I (Security Properties): Secure by default, protection from unauthorised access, data confidentiality
%      and integrity, least privilege, encrypted communications, minimal attack surface.
%    - Part II (Vulnerability Handling): Maintain an automated, machine-readable Software Bill of Materials
%      (SBOM) in CycloneDX/SPDX; establish Coordinated Vulnerability Disclosure (security.txt); deliver security-only
%      patches separated from feature updates; define support period commitment.
%
% 3. Compliance Deadlines and Reporting Obligations:
%    - 11 September 2026: Mandatory reporting of actively exploited vulnerabilities and severe incidents
%      within 24 hours (early warning) to CSIRT / CNCS and ENISA.
%    - 11 December 2027: Full application of CRA requirements (conformity assessment, CE marking, EU declaration
%      of conformity, technical documentation file).
%
% 4. National Supervision:
%    - Supervision and market surveillance authority in Portugal: Centro Nacional de Cibersegurança (CNCS).
